% Options for packages loaded elsewhere
\PassOptionsToPackage{unicode}{hyperref}
\PassOptionsToPackage{hyphens}{url}
%
\documentclass[
  ignorenonframetext, 10pt
]{beamer}
\usepackage{pgfpages}
\setbeamertemplate{caption}[numbered]
\setbeamertemplate{caption label separator}{: }
\setbeamercolor{caption name}{fg=normal text.fg}
\beamertemplatenavigationsymbolsempty
% Prevent slide breaks in the middle of a paragraph
\widowpenalties 1 10000
\raggedbottom
\setbeamertemplate{part page}{
  \centering
  \begin{beamercolorbox}[sep=16pt,center]{part title}
    \usebeamerfont{part title}\insertpart\par
  \end{beamercolorbox}
}
\setbeamertemplate{section page}{
  \centering
  \begin{beamercolorbox}[sep=12pt,center]{part title}
    \usebeamerfont{section title}\insertsection\par
  \end{beamercolorbox}
}
\setbeamertemplate{subsection page}{
  \centering
  \begin{beamercolorbox}[sep=8pt,center]{part title}
    \usebeamerfont{subsection title}\insertsubsection\par
  \end{beamercolorbox}
}
\AtBeginPart{
  \frame{\partpage}
}
\AtBeginSection{
  \ifbibliography
  \else
    \frame{\sectionpage}
  \fi
}
\AtBeginSubsection{
  \frame{\subsectionpage}
}
\usepackage{amsmath,amssymb}
\usepackage{iftex}
\ifPDFTeX
  \usepackage[T1]{fontenc}
  \usepackage[utf8]{inputenc}
  \usepackage{textcomp} % provide euro and other symbols
\else % if luatex or xetex
  \usepackage{unicode-math} % this also loads fontspec
  \defaultfontfeatures{Scale=MatchLowercase}
  \defaultfontfeatures[\rmfamily]{Ligatures=TeX,Scale=1}
\fi
\usepackage{lmodern}
\ifPDFTeX\else
  % xetex/luatex font selection
\fi
% Use upquote if available, for straight quotes in verbatim environments
\IfFileExists{upquote.sty}{\usepackage{upquote}}{}
\IfFileExists{microtype.sty}{% use microtype if available
  \usepackage[]{microtype}
  \UseMicrotypeSet[protrusion]{basicmath} % disable protrusion for tt fonts
}{}
\makeatletter
\@ifundefined{KOMAClassName}{% if non-KOMA class
  \IfFileExists{parskip.sty}{%
    \usepackage{parskip}
  }{% else
    \setlength{\parindent}{0pt}
    \setlength{\parskip}{6pt plus 2pt minus 1pt}}
}{% if KOMA class
  \KOMAoptions{parskip=half}}
\makeatother
\usepackage{xcolor}
\newif\ifbibliography
\setlength{\emergencystretch}{3em} % prevent overfull lines
\providecommand{\tightlist}{%
  \setlength{\itemsep}{0pt}\setlength{\parskip}{0pt}}
\setcounter{secnumdepth}{-\maxdimen} % remove section numbering
%\usetheme{metropolis}


%\setbeamercolor{normal text}{bg=white}
%\addtobeamertemplate{frametitle}{}{\vspace{-0.5em}} % decrease
%\setbeamersize{text margin left=5mm,text margin right=10mm}
%\usefonttheme[onlymath]{serif}
%\usefonttheme{professionalfonts}
%%%%%%%%%%%%%%%%
% Packages
%%%%%%%%%%%%%%%%

\usepackage{geometry}
\usepackage{graphicx}
\usepackage{amssymb}
\usepackage{amsopn}
\usepackage{cancel}
\usepackage{epstopdf}
\usepackage{amsmath}  	% this permits text in eqnarray among other benefits
\usepackage{url}		% produces hyperlinks
\usepackage{hyperref}	% allows for color usage in tables
\usepackage[english]{babel}
%\usepackage[latin1]{inputenc}
\usepackage{colortbl}	% allows for color usage in tables
\usepackage{multirow}	% allows for rows that span multiple rows in tables
\usepackage{color}		% this package has a variety of color options
\usepackage{xcolor}
\usepackage{pgf}
\usepackage{calc}
\usepackage{ulem}
\usepackage{multicol}
\usepackage{booktabs}
\usepackage{textcomp}
%\usepackage[misc]{ifsym} % for the dice symbol \Cube{}
%\usepackage{txfonts}
%\usepackage{listings}
%\usepackage{tikz}
\usepackage{array}
\usepackage{wasysym}
\usepackage{fancyvrb}

%% change fontsize of R code
\ifdefined\Shaded
\let\oldShaded\Shaded
\let\endoldShaded\endShaded
\renewenvironment{Shaded}{\footnotesize\oldShaded}{\endoldShaded}
\fi
%% change fontsize of output
\let\oldverbatim\verbatim
\let\endoldverbatim\endverbatim
\renewenvironment{verbatim}{\footnotesize\oldverbatim}{\endoldverbatim}

%%%%%%%%%%%%%%%%
% Remove navigation symbols
%%%%%%%%%%%%%%%%

\setbeamertemplate{navigation symbols}{}

\setbeamertemplate{footline}[frame number]

%%%%%%%%%%%%%%%%
% User defined colors
%%%%%%%%%%%%%%%%

\xdefinecolor{oiB}{rgb}{0.22,0.52,0.72}
\definecolor{oiG}{rgb}{.298,.447,.114}
\xdefinecolor{hlblue}{rgb}{0.051,0.65,1}
\xdefinecolor{gray}{rgb}{0.5, 0.5, 0.5}
\xdefinecolor{darkGray}{rgb}{0.3, 0.3, 0.3}
\xdefinecolor{darkerGray}{rgb}{0.2, 0.2, 0.2}
\xdefinecolor{rubineRed}{rgb}{0.89,0,0.30}
\xdefinecolor{irishGreen}{rgb}{0,0.60,0}	
\definecolor{lightGreen}{rgb}{0.387,0.581,0.148}

%%%%%%%%%%%%%%%%
% Template colors
%%%%%%%%%%%%%%%%

%\setbeamercolor*{palette primary}{fg=white,bg= oiB!80!black!90}
%\setbeamercolor*{palette secondary}{fg=black,bg= oiB!80!black}
%\setbeamercolor*{palette tertiary}{fg=white,bg= oiB!80!black!80}
%\setbeamercolor*{palette quaternary}{fg=white,bg= oiB}
%\setbeamercolor{structure}{fg= oiB}
%\setbeamercolor{frametitle}{bg= oiB!90}
%\setbeamertemplate{blocks}[shadow=false]
%\setbeamersize{text margin left=2em,text margin right=2em}

%\setbeamercolor{code body}{bg=gray!20!white!80,fg=black}


%%%%%%%%%%%%%%%%
% Get rid of fancy enumerated list bullets
%%%%%%%%%%%%%%%%

%\setbeamertemplate{enumerate items}[default]

%%%%%%%%%%%%%%%%
% Custom commands
%%%%%%%%%%%%%%%%

\newcommand{\Sum}{\sum\nolimits}
\newcommand{\Hide}[2]{\underline{~~~\onslide<#1>{\alert{#2}}~~~}}
\newcommand{\hide}[2]{\underline{~~\onslide<#1>{\alert{#2}}~~}}
\newcommand{\hid}[2]{\onslide<#1>{\alert{#2}}}
\newcommand{\code}[1]{\textcolor{blue}{\texttt{#1}}}

% red
\newcommand{\red}[1]{\textit{\textcolor{rubineRed}{#1}}}

% pink
\newcommand{\pink}[1]{\textit{\textcolor{rubineRed!90!white!50}{#1}}}

% green
\newcommand{\green}[1]{\textit{\textcolor{irishGreen}{#1}}}

% orange
\newcommand{\orange}[1]{\textit{\textcolor{orange}{#1}}}

% links: webURL, webLin, appLink
\newcommand{\webURL}[1]{\urlstyle{same}{ \textit{\textcolor{darkGray}{\url{#1}}}}}
\newcommand{\webLink}[2]{\href{#1}{\textcolor{darkGray}{{#2}}}}
\newcommand{\appLink}[2]{\href{#1}{\textcolor{white}{{#2}}}}

% mail
\newcommand{\mail}[1]{\href{mailto:#1}{\textit{\textcolor{darkGray}{#1}}}}

% highlighting: hl, hlGr, mathhl
\newcommand{\hl}[1]{\textit{\textcolor{hlblue}{#1}}}
\newcommand{\hlGr}[1]{\textit{\textcolor{lightGreen}{#1}}}
\newcommand{\mathhl}[1]{\textcolor{hlblue}{\ensuremath{#1}}}

% two col: two columns
\newenvironment{twocol}[4]{
\begin{columns}[c]
\column{#1\textwidth}
#3
\column{#2\textwidth}
#4
\end{columns}
}





%%%%%%%%%%%%%%%%
% Change margin
%%%%%%%%%%%%%%%%

\newenvironment{changemargin}[2]{%
\begin{list}{}{%
\setlength{\topsep}{0pt}%
\setlength{\leftmargin}{#1}%
\setlength{\rightmargin}{#2}%
\setlength{\listparindent}{\parindent}%
\setlength{\itemindent}{\parindent}%
\setlength{\parsep}{\parskip}%
}%
\item}{\end{list}}

%%%%%%%%%%%%%%%%
% Footnote
%%%%%%%%%%%%%%%%

\long\def\symbolfootnote[#1]#2{\begingroup%
\def\thefootnote{\fnsymbol{footnote}}\footnote[#1]{#2}\endgroup}

%%%%%%%%%%%%%%%%
% Graphics
%%%%%%%%%%%%%%%%

\DeclareGraphicsRule{.tif}{png}{.png}{`convert #1 `dirname #1`/`basename #1 .tif`.png}



\setbeamercolor{normal text}{bg=white}

\DeclareMathOperator{\var}{var}
\newcommand{\boxmodel}[1]{\begin{array}{|c|}#1\\[2pt]\hline\end{array}}
\newcommand{\BXX}[2]{\begin{array}{|@{\,}c@{\,}|@{\,}c@{\,}|}\hline#1&#2\\[-1pt]\hline\end{array}}
\newcommand{\p}{\mathrm{P}}
\newcommand{\SE}{\mathrm{SE}}
\newcommand{\xbar}{\overline{x}}
\newcommand{\Xbar}{\overline{X}}
\newcommand{\ybar}{\overline{y}}
\newcommand{\Ybar}{\overline{Y}}
\newcommand{\zbar}{\overline{z}}
\newcommand{\Zbar}{\overline{Z}}
\newcommand{\wbar}{\overline{w}}
\newcommand{\Wbar}{\overline{W}}
\newcommand{\vbar}{\overline{v}}
\newcommand{\Vbar}{\overline{V}}
\newcommand{\dbar}{\overline{d}}
\newcommand{\Dbar}{\overline{D}}
\newcommand{\ebar}{\overline{e}}
\newcommand{\ehat}{\widehat{\epsilon}}
\newcommand{\yhat}{\widehat{y}}
\newcommand{\Yhat}{\widehat{Y}}
\newcommand{\bhat}{\widehat{\beta}}
\newcommand{\dl}[1]{\widehat{\delta}_{#1}}
\newcommand{\g}[1]{\widehat{\gamma}_{#1}}
\newcommand{\al}[1]{\widehat{\alpha}_{#1}}
\newcommand{\Bj}{{\beta_j}}
\newcommand{\Bp}{{\beta_p}}
\newcommand{\Bq}{{\beta_q}}
\renewcommand{\b}[1]{\widehat{\beta}_{#1}}
\newcommand{\B}[1]{{\beta}_{#1}}
\newcommand{\BETA}{{\boldsymbol\beta}}
\newcommand{\bp}{\widehat{\beta}_p}
\newcommand{\bq}{\widehat{\beta}_q}
\newcommand{\bj}{\widehat{\beta}_j}
\newcommand{\df}{\mathrm{df}}
\DeclareMathOperator{\VIF}{VIF}
\DeclareMathOperator{\SSE}{SSE}
\DeclareMathOperator{\SST}{SST}
\DeclareMathOperator{\SSR}{SSR}
\DeclareMathOperator{\MSE}{MSE}
\DeclareMathOperator{\MST}{MST}
\DeclareMathOperator{\MSR}{MSR}
\DeclareMathOperator{\RM}{RM}
\DeclareMathOperator{\FM}{FM}
\DeclareMathOperator{\E}{E}
\DeclareMathOperator{\SD}{SD}
\DeclareMathOperator{\se}{se}
\DeclareMathOperator{\cov}{Cov}
\DeclareMathOperator{\corr}{Corr}
\DeclareMathOperator{\cor}{Cor}
\DeclareMathOperator{\Var}{Var}
\DeclareMathOperator{\logit}{logit}
\newcommand{\hsigma}{\widehat{\sigma}}
\newcommand{\betahat}{\widehat{\beta}}
\newcommand{\slide}{<br><br><br><br><br>}
\newcommand{\bbeta}{\boldsymbol{\beta}}
\newcommand{\bX}{\mathbf{X}}
\newcommand{\bY}{\mathbf{Y}}
\newcommand{\bH}{\mathbf{H}}
\newcommand{\bI}{\mathbf{I}}
\newcommand{\be}{\mathbf{e}}
\newcommand{\pval}{\mathit{P}\mathrm{-val}}



\newcommand{\vmu}{\mbox{\boldmath{$\mu$}}}
\newcommand{\vR}{\mbox{\boldmath{$R$}}}
\newcommand{\vpsi}{\mbox{\boldmath{$\psi$}}}
\newcommand{\vZ}{\mbox{\boldmath{$Z$}}}
\newcommand{\vX}{\mbox{\boldmath{$X$}}}
\newcommand{\vY}{\mbox{\boldmath{$Y$}}}
\newcommand{\vW}{\mbox{\boldmath{$W$}}}
\newcommand{\vz}{\mbox{\boldmath{$z$}}}
\newcommand{\vr}{\mbox{\boldmath{$r$}}}
\newcommand{\vx}{\mbox{\boldmath{$x$}}}
\newcommand{\vk}{\mbox{\boldmath{$k$}}}
\newcommand{\vp}{\mbox{\boldmath{$p$}}}
\newcommand{\vv}{\mbox{\boldmath{$v$}}}
\newcommand{\vg}{\mbox{\boldmath{$g$}}}
\newcommand{\vbeta}{\mbox{\boldmath{$\beta$}}}
\newcommand{\vone}{\mbox{\boldmath{$1$}}}
\newcommand{\vep}{\mbox{\boldmath{$\epsilon$}}}
\newcommand{\veta}{\mbox{\boldmath{$\eta$}}}
\newcommand{\vl}{\mbox{\boldmath{$l$}}}
\newcommand{\ve}{\mbox{\boldmath{$e$}}}
\newcommand{\va}{\mbox{\boldmath{$a$}}}
\newcommand{\vb}{\mbox{\boldmath{$b$}}}
\newcommand{\vPsi}{\mbox{\boldmath{$\Psi$}}}
\DeclareMathOperator*{\V}{\mathbb{V}}


% geometry: topsep=0pt 
\ifLuaTeX
  \usepackage{selnolig}  % disable illegal ligatures
\fi
\IfFileExists{bookmark.sty}{\usepackage{bookmark}}{\usepackage{hyperref}}
\IfFileExists{xurl.sty}{\usepackage{xurl}}{} % add URL line breaks if available
\urlstyle{same}
\hypersetup{
  pdftitle={STAT 20000 Class Meetings 2 \& 3Chapter 1-2 Experiments \& Observaitonal Studies},
  hidelinks,
  pdfcreator={LaTeX via pandoc}}

\title{STAT 20000 Class Meetings 2 \& 3\newline Chapter 1-2 Experiments
\& Observaitonal Studies}
\author{}
\date{\vspace{-2.5em}}

\begin{document}
\frame{\titlepage}

\begin{frame}{Explanatory Variable \& Response}
\protect\hypertarget{explanatory-variable-response}{}
For each of the following claims, please identify the explanatory
variable and the response

\begin{enumerate}
\tightlist
\item
  ``\textit{Daily Exercise Improves Mental Performance}''

  \begin{itemize}
  \tightlist
  \item
    Explanatory: Daily Exercise
  \item
    Response: Mental Health
  \end{itemize}
\item
  ``\textit{Want to lose weight? Watch more TV shows!}''

  \begin{itemize}
  \tightlist
  \item
    Explanatory: Watching More TV
  \item
    Response: Weight Loss
  \end{itemize}
\item
  ``\textit{More ice creams sales have been linked to more deaths by drowning}''

  \begin{itemize}
  \tightlist
  \item
    Ice Cream Sales
  \item
    Count of Death by Drowning
  \end{itemize}
\end{enumerate}
\end{frame}

\begin{frame}{Association and Causation}
\protect\hypertarget{association-and-causation}{}
Two variables are \textbf{associated} if values of one variable tend to
be related to values of the other variable. \bigskip

Two variables have a \textbf{causal (or ``cause-and-effect'') relation},
if changing the value of the explanatory variable influences the value
of the response variable
\end{frame}

\begin{frame}{Should We Buy More TVs to Live Longer?}
\protect\hypertarget{should-we-buy-more-tvs-to-live-longer}{}
\vspace{-8pt}

\begin{center}\includegraphics[width=0.8\linewidth]{ch02_24spr_files/figure-beamer/unnamed-chunk-1-1} \end{center}

\vspace{-12pt}
\Large

Association does not imply causation! \scriptsize

Source: \url{http://jse.amstat.org/v2n2/datasets.rossman.html} \newline
``Televisions, Physicians, and Life Expectancy'' by Allan J.
Rossman\newline \textit{Journal of Statistics Education} v.2, n.2 (1994)
\end{frame}

\begin{frame}{Confounding Variable}
\protect\hypertarget{confounding-variable}{}
\textbf{Definition}: A third variable associated with both the
explanatory variable and the response is called a
\hl{confounding variable}

\begin{itemize}
\item
  A confounding variable can offer a plausible explanation for an
  association between the explanatory and the response variables
\item
  \red{Whenever confounding variables are (potentially) present,
  a causal relationship  cannot be established}
\end{itemize}
\end{frame}

\begin{frame}{Confounding Variable}
\protect\hypertarget{confounding-variable-1}{}
\includegraphics[width=1\linewidth]{ch02_24spr_files/figure-beamer/unnamed-chunk-2-1}
\end{frame}

\begin{frame}{TVs and Life Expectancy}
\protect\hypertarget{tvs-and-life-expectancy}{}
\includegraphics[width=1\linewidth]{ch02_24spr_files/figure-beamer/unnamed-chunk-3-1}
\end{frame}

\begin{frame}{Identify a Confounding Variable (1)}
\protect\hypertarget{identify-a-confounding-variable-1}{}
Association: More ice creams sales have been linked to more deaths by
drowning

\includegraphics[width=1\linewidth]{ch02_24spr_files/figure-beamer/unnamed-chunk-4-1}
\end{frame}

\begin{frame}{Identify a Confounding Variable (2)}
\protect\hypertarget{identify-a-confounding-variable-2}{}
Association: The total amount of beef consumed and the total amount of
pork consumed worldwide are positively related over the past 100 years

\includegraphics[width=1\linewidth]{ch02_24spr_files/figure-beamer/unnamed-chunk-5-1}
\end{frame}

\begin{frame}{Identify a Confounding Variable (3)}
\protect\hypertarget{identify-a-confounding-variable-3}{}
Association: People with shorter hair tend to be taller.

\includegraphics[width=1\linewidth]{ch02_24spr_files/figure-beamer/unnamed-chunk-6-1}
\end{frame}

\begin{frame}{Experiment v.s. Observational Study}
\protect\hypertarget{experiment-v.s.-observational-study}{}
An \textbf{observational study} is a study in which the researcher does
not actively control the value of any variable, but simply observes the
value as they naturally exist.\bigskip

An \textbf{experiment} is a study in which the researcher actively
controls one or more of the explanatory variables
\end{frame}

\begin{frame}{Observational Studies}
\protect\hypertarget{observational-studies}{}
\begin{itemize}
\tightlist
\item
  Many studies in social science and public health are observational
  studies.
\item
  Example: to study the effect of smoking on health

  \begin{itemize}
  \tightlist
  \item
    Treatment group (``exposed group''): smokers;
  \item
    Control group: no-smokers;
  \end{itemize}
\item
  Since the investigators have no control over allocation, the treatment
  group and control group might be different in important aspects other
  than the treatment, i.e.,\newline \textbf{there are many confounding
  variables}.
\end{itemize}

Observational studies only established \textbf{association}, NOT
\textbf{causation}!
\end{frame}

\begin{frame}{Randomization}
\protect\hypertarget{randomization}{}
How can we make sure to avoid confounding variables?

\begin{center}
\renewcommand{\arraystretch}{1.6}
\begin{tabular}{|c|}\hline
\LARGE\red{RANDOMLY} assign\\
\LARGE values of the\\
\LARGE explanatory variable\\\hline
\end{tabular}
\end{center}
\pause

\begin{itemize}
\tightlist
\item
  For experiments,the explanatory variable is ``treated or not''
\item
  Randomly \(\neq\) Haphazardly
\item
  Avoid the human factor: use coin tossing/random number generation.
\end{itemize}
\end{frame}

\begin{frame}{Controlled Experiment on Clofibrate (FPP, Ch 2, Sec 2)}
\protect\hypertarget{controlled-experiment-on-clofibrate-fpp-ch-2-sec-2}{}
\normalsize
\begin{center}
\begin{tabular}{lcc@{~~~}cc}
& \multicolumn{2}{c}{Clofibrate} & \multicolumn{2}{c}{Placebo}\\
\midrule
 & Number & Deaths & Number & Deaths\\

\only<2->{Adherers     & ~708 & 15\% & \only<3->{1813 & 15\%}}\\
\only<2->{Non-adherers & ~357 & 25\% & \only<3->{~882 & 28\%}}\\
\midrule
Total group    & 1103 & 20\% & 2798 & 21\%
\end{tabular}
\end{center}
\large
\begin{itemize}
\item<1-> Is Clofibrate effective?
\item<2-> Adherers v.s. non-adherers in the treatment group\\
      Does this mean that Clofibrate is actually effective?\bigskip\bigskip
%\item Adhering or not is {\bf not} controlled by the
%investigators.  So the comparison is {\bf observational}, and does not support the claim that the drug is effective.
\item<3-> Adherers v.s. non-adherers in the control group
\item<4-> Conclusion
\begin{itemize}\large
\item Clofibrate has no effect;
\item Adherers are different from non-adherers;
\end{itemize}
\end{itemize}
\end{frame}

\begin{frame}{Incentive for Smoke Quitting}
\protect\hypertarget{incentive-for-smoke-quitting}{}
Smokers at a big company were invited to participate in a smoking
cessation program and randomly assigned to one of two groups.

\begin{itemize}
\tightlist
\item
  \textbf{Reward Only}: subjects got a cash award if stopped smoking for
  6 months
\item
  \textbf{Deposit}: subjects had to deposit some money first, could get
  deposit back + a bonus if stopped smoking.
\end{itemize}

\begin{center}
\begin{tabular}{@{}r|cc|r}
\multicolumn{1}{c}{}&\multicolumn{2}{r}{\textbf{\large Reward Only Group}} \\
\multicolumn{1}{c}{}& Quit Smoking & Not Quit & Total\\\hline
  Agreed to Participate & 156 & 758 & 914\\
Declined to Participate &   3 & 100 & 103\\\hline
Total & 159 & 858 & 1017
\end{tabular}\\[5pt]
\begin{tabular}{@{}r|cc|r@{}}
\multicolumn{1}{c}{}&\multicolumn{2}{c}{\textbf{\large Deposit Group}}& \\
\multicolumn{1}{c}{}& Quit Smoking & Not Quit & Total\\\hline
  Agreed to Participate & 78 &  68 & 146\\
Declined to Participate & 30 & 877 & 907\\\hline
Total & 108 & 945 & 1053
\end{tabular}
\end{center}
\end{frame}

\begin{frame}{}
\protect\hypertarget{section}{}
Which of the following is the correct and best conclusion of this
experiment? Explain your choice.

\begin{enumerate}
[(a)]
\tightlist
\item
  ``Reward'' is more effective than ``Deposit'' since the success rate
  of the entire Reward-Only Group \(159/1017\approx15.6\)\% is higher
  than that of the entire Deposit Group \(108/1053\approx10.3\)\% and
  the two groups are divided by randomization.
\item
  ``Deposit'' seems more effective than ``Deposit'' since the success
  rates among those to agree to participate in the Deposit Group
  \(78/146\approx 53.4\%\) is higher than that for the Reward Group,
  \(156/914\approx 17.1\%\). However, this comparison cannot establish
  causality since people self-selected whether to participate in the
  assigned program.
\item
  Same comparison as (b) but we CAN make a causal conclusion that
  ``Deposit'' is better since this is a randomized experiment
\item
  Same comparison as (a) but we CANNOT make a causal conclusion that
  ``Reward'' is better due to the huge difference between the
  participation rates of the two groups.
\end{enumerate}
\end{frame}

\end{document}
